\documentclass[11pt]{article}
\usepackage[margin=1in]{geometry}
\usepackage{amsmath,amssymb}
\title{A cluster-tilted algebra of representation dimension zero}
\author{The Last Math Competition, conjecture 00000004041}
\date{}
\begin{document}
\maketitle

Let \(k\) be a field and use the hereditary algebra of the type \(A_1\)
quiver, which is \(k\). Its cluster category is the orbit category
\(\mathcal C=D^b(\mathrm{mod}\,k)/\langle\tau^{-1}[1]\rangle\)
\cite{bmrtt,bmr}. The Serre functor of \(D^b(\mathrm{mod}\,k)\) is the
identity, so \(\tau=[-1]\) and the orbit functor is \([2]\). Every bounded
complex of finite-dimensional \(k\)-vector spaces is isomorphic in the
derived category to the direct sum of its shifted cohomology. Thus every
object of \(\mathcal C\) is a sum of copies of \(S_0=k\) and \(S_1=k[1]\).
Indeed, the orbit morphism space is
\[
  \operatorname{Hom}_{\mathcal C}(S_i,S_j)
   =\bigoplus_{m\in\mathbb Z}\operatorname{Hom}_{D^b(k)}(k[i],k[j+2m]) .
\]
Because \(k\) is semisimple, there is a one-dimensional term exactly when
\(i=j+2m\). For \(i,j\in\{0,1\}\), the morphism space is therefore \(k\) if
\(i=j\) and zero otherwise. The shift \([1]\) interchanges the two objects.
Consequently, for \(T=S_0\), the condition
\(\operatorname{Ext}^1_{\mathcal C}(T,X)=0\) holds exactly when \(X\) is a
sum of copies of \(S_0\). Hence \(S_0\) is cluster-tilting, as is \(S_1\).
The zero object is not cluster-tilting and the sum of both simples is not
rigid. There are exactly two basic cluster-tilting objects, joined by the
mutation exchanging the complements of the zero almost complete object.
In particular the mutation graph is finite. At the quiver level, the
type \(A_1\) quiver also has just one mutation class representative, the
single vertex with no arrows.

By the definition of a cluster-tilted algebra \cite{bmr},
\(A=\operatorname{End}_{\mathcal C}(T)\cong k\) is cluster-tilted. The
identification respects multiplication because the endomorphism space
has only its degree-zero term and composition there is multiplication
of scalars. Use the representation dimension given by the infimum of
\(\operatorname{gl.dim}\operatorname{End}_A(M)\) over generator-cogenerators
\(M\) \cite{xi}. For \(A=k\), the module \(M=k\) is a
generator-cogenerator since every finite-dimensional vector space is a
finite sum of copies of it. Its endomorphism ring is \(k\), whose global
dimension is zero because every vector space is free and projective.
Global dimensions are nonnegative, so \(\operatorname{repdim}(A)=0\).
This contradicts the claimed restriction to representation dimension
two or three in the finite mutation case.

\paragraph{Lean verification.} The Lean model records the multiplicities
of the two shift classes. It verifies their orbit Hom dimensions, the
shifted Ext dimensions, the complete list of basic cluster-tilting
objects, and the exchange graph. Its final statement combines the finite
mutation data with the zero representation dimension computed above.

\begin{thebibliography}{3}
\bibitem{bmrtt} A.~B. Buan, R.~J. Marsh, M. Reineke, I. Reiten and
G. Todorov, \emph{Tilting theory and cluster combinatorics}, Adv. Math.
\textbf{204} (2006), 572--618.
\bibitem{bmr} A.~B. Buan, R.~J. Marsh and I. Reiten,
\emph{Cluster-tilted algebras}, Trans. Amer. Math. Soc. \textbf{359}
(2007), 323--332.
\bibitem{xi} C. Xi, \emph{Representation dimension and quasi-hereditary
algebras}, Adv. Math. \textbf{168} (2002), 193--212.
\end{thebibliography}
\end{document}
